\section{Multi-step autoregressive ``staircase'' diagnostic}\label{sec:cmp_intflux}

Single-step CSI evaluates the model on the next 4 frames only.
For mission-relevant lead times we need the model to chain its own
predictions. We instantiated a staircase protocol on HARP 11930
(\emph{train} cube, qualitative): input \(T_{\text{in}}=10\) real
frames, predict \(T_{\text{out}}=4\), commit the first
\texttt{stride}\(=2\) predicted frames into the input buffer, slide,
repeat for 20 steps (i.e.\ 40 committed frames \(\equiv\) 8\,h forecast
lead). Full-frame inference at each step (tiled, stride 64, no
overlap). Pure autoregressive --- the model never sees a real frame
past \(T_{\text{in}}\) after step 0.

We ran the protocol on S3 (best single-step regression on this cube),
S11 (best single-step regression on the test cubes), and S16 (new
SOTA).

\subsection*{Headline numbers}
\begin{center}\small
\begin{tabular}{l r r r}
\toprule
Arm & MAE (integrated flux) & Var.\ ratio & Drift corr.\ \(\rho\)\\
\midrule
S3  & \(1.5\!\times\!10^8\) & 0.028 (mode collapse)      & +0.14 \\
S11 & \(7.6\!\times\!10^8\) & 15.8 (runaway explosion)   & +0.81 \\
S16 & (see below)                 & ---                       & ---   \\
\bottomrule
\end{tabular}
\end{center}

\emph{(S16 staircase numbers are in \texttt{outputs/staircase\_s16\_harp\_11930/series.npz}; integrated-flux plot below.)}

\subsection*{S3: mode collapse}
Within 4 chained frames the predicted field amplitude drops from
\(|\hat w|_\infty\!=\!1.3\!\times\!10^7\) to
\(\sim\!10^5\) (100\(\times\) collapse). Frame-to-frame field
difference settles to \(\sim\)150 on values of order \(1.5\!\times\!10^5\)
(\(<\)0.1\% relative motion). After collapse the network produces a
\emph{static} low-amplitude field for the remaining 36 frames. MAE
stays low only because the static value happens to be close to the
GT temporal mean.

\noindent\includegraphics[width=0.49\linewidth]{s3_staircase_4steps}
\includegraphics[width=0.49\linewidth]{s3_staircase_intflux}

\subsection*{S11: runaway explosion}
Behaves like S3 for the first \(\sim\)9 chained steps (low-amplitude
collapse phase) and then \emph{variance explodes}: frame
\(|\hat w|_\infty\) grows from \(\sim\!10^5\) to
\(2\!\times\!10^6\) by frame 39; per-frame field difference grows from
150 to \(\sim\)3000. The 12-step grid below makes the transition
visible --- the bottom rows are progressively more saturated.
Variance ratio 15.8\(\times\), drift correlation +0.81 (error grows
strongly with horizon).

\noindent\includegraphics[width=0.49\linewidth]{s11_staircase_4steps}
\includegraphics[width=0.49\linewidth]{s11_staircase_intflux}

\noindent\includegraphics[width=\linewidth]{s11_staircase_12steps}

\subsection*{S16: staircase}
\noindent\includegraphics[width=0.49\linewidth]{s16_staircase_4steps}
\includegraphics[width=0.49\linewidth]{s16_staircase_intflux}

\noindent\includegraphics[width=\linewidth]{s16_staircase_12steps}

\subsection*{Verdict}
None of S3 / S11 sustains chained prediction --- one collapses, the
other explodes. The two failure modes are symmetric: L1-dominated
losses converge to a stable low-amplitude fixed point, while
amplitude-rewarding losses (S11's extreme weight + fast-TF) shift
the autoregressive map to an unstable regime. Long-horizon
forecasting is genuinely open; the single-step gains from S16 should
be re-tested under chained rollouts before claiming
multi-hour skill.
